\documentclass[10pt,letterpaper]{article}
\usepackage[T1]{fontenc}
\usepackage[utf8]{inputenc}
\usepackage[sfdefault]{sourcesanspro}
\usepackage{inconsolata}
\usepackage[margin=0.84in,headheight=15pt]{geometry}
\usepackage{microtype}
\usepackage{xcolor}
\usepackage{booktabs,longtable,tabularx,array}
\usepackage{enumitem}
\usepackage{fancyhdr,lastpage}
\usepackage[most]{tcolorbox}
\usepackage{listings}
\usepackage{titlesec}
\usepackage{hyperref}
\usepackage{amsmath,amssymb}
\usepackage{multicol}
\definecolor{navy}{HTML}{15324C}
\definecolor{teal}{HTML}{087E83}
\definecolor{slate}{HTML}{536879}
\definecolor{paper}{HTML}{F2F6F8}
\definecolor{warm}{HTML}{E8F4F3}
\hypersetup{colorlinks=true,urlcolor=teal,linkcolor=navy,pdfauthor={Software Engineering Reference Suite},pdfsubject={MISRA C++:2023 engineering guidance}}
\titleformat{\section}{\Large\bfseries\color{navy}}{\thesection}{0.7em}{}
\titleformat{\subsection}{\large\bfseries\color{teal}}{\thesubsection}{0.6em}{}
\titleformat{\subsubsection}{\normalsize\bfseries\color{navy}}{\thesubsubsection}{0.5em}{}
\setlist[itemize]{leftmargin=*,topsep=2pt,itemsep=2pt}
\setlist[enumerate]{leftmargin=*,topsep=2pt,itemsep=3pt}
\setlength{\parskip}{5pt}
\setlength{\parindent}{0pt}
\setlength{\emergencystretch}{2em}
\pagestyle{fancy}
\fancyhf{}
\fancyhead[L]{\small\color{slate}MISRA C++:2023 Engineering Suite}
\fancyhead[R]{\small\color{slate}\nouppercase{\leftmark}}
\fancyfoot[L]{\scriptsize\color{slate}Independent engineering guidance; see licensed source}
\fancyfoot[R]{\scriptsize\color{slate}\thepage\ / \pageref{LastPage}}
\newtcolorbox{sourcebox}{colback=paper,colframe=navy!30,boxrule=0.45pt,arc=1.6pt,breakable}
\newtcolorbox{practicebox}{colback=warm,colframe=teal!40,boxrule=0.45pt,arc=1.6pt,breakable}
\newcommand{\SourceNote}[1]{\begin{sourcebox}\textbf{SOURCE BASIS / NORMATIVE POINTER}\quad #1\end{sourcebox}}
\newcommand{\PracticeNote}[1]{\begin{practicebox}\textbf{ENGINEERING RECOMMENDATION / NOT MISRA TEXT}\quad #1\end{practicebox}}
\newcommand{\Record}[1]{\textsf{\bfseries #1}}
\lstdefinestyle{cppsuite}{language=C++,basicstyle=\ttfamily\footnotesize,keywordstyle=\color{navy}\bfseries,commentstyle=\color{slate},stringstyle=\color{teal},breaklines=true,breakatwhitespace=true,columns=fullflexible,frame=single,framerule=0.3pt,rulecolor=\color{navy!35},backgroundcolor=\color{paper},showstringspaces=false,tabsize=2,xleftmargin=0.5em,xrightmargin=0.5em,aboveskip=4pt,belowskip=4pt}
\lstset{style=cppsuite}
\newcommand{\SuiteTitle}[3]{%
\begin{titlepage}\thispagestyle{empty}
\vspace*{1.5cm}
{\color{teal}\rule{2.4cm}{2pt}}\par\vspace{1cm}
{\Huge\bfseries\color{navy}MISRA C++:2023\\[0.4em]#1}\par\vspace{0.7cm}
{\Large\color{slate}#2}\par\vspace{1.2cm}
\begin{sourcebox}\textbf{Source baseline:} MISRA C++:2023, Guidelines for the use of C++17 in critical systems, October 2023 (user-supplied licensed copy). This is an independent companion, not an official MISRA publication, rule replacement, certified checker, or claim of project compliance.\end{sourcebox}
\vfill{\large\color{navy}Technical edition \textbullet\ 16 October 2026}\par\smallskip
{\small\color{slate}#3}\par\end{titlepage}
\tableofcontents\clearpage
}
\newcommand{\StandardsBoundary}{%
\section*{How to read normative status}\addcontentsline{toc}{section}{How to read normative status}
\textbf{Normative source:} the licensed MISRA C++:2023 guideline text, including any applicable \emph{Amplification} and \emph{Exception}. A guideline \emph{headline} alone is not a complete compliance test. The source's Section~3 defines classifications, categories, decidability, analysis scope, and applicability.\par
\textbf{Independent guidance:} all review questions, programming patterns, lifecycle processes, sample thresholds, ownership assignments, controls, and CI configurations in this suite are \emph{engineering recommendations}, not statements issued by MISRA.\par
\textbf{Critical limitation:} no individual example is asserted compliant with the whole standard. For formal compliance claims, follow the applicable edition of \emph{MISRA Compliance} as the source itself instructs, with a documented compliance plan and supporting evidence.\par
}

\begin{document}
\SuiteTitle{Static Analysis and Tool Assurance}{Rule coverage, analyzer limits, technical qualification and triage}{Volume 06 / Verification}
\StandardsBoundary

\section{Verification architecture}
Use multiple independent evidence channels: compiler diagnostics, rule-aware static analysis, cross-translation-unit data-flow analysis, human review, unit/integration testing, and runtime checks in testing environments. No single one of these substitutes for the whole compliance process.

The source states that directives cannot be decided using source alone; rules may be decidable or undecidable and may require single-translation-unit or system-level analysis. The goal is to demonstrate coverage and limitation handling, not maximize an arbitrary warning count.


\section{Tool selection and configuration qualification}
\small\setlength{\tabcolsep}{3pt}\renewcommand{\arraystretch}{1.15}
\begin{longtable}{p{0.18\linewidth}p{0.41\linewidth}p{0.33\linewidth}}
\toprule
Evaluation axis & Qualification question & Evidence\\ \midrule
\endfirsthead
\toprule
Evaluation axis & Qualification question & Evidence\\ \midrule
\endhead
C++17 dialect & Does parser faithfully model target C++17, vendor extensions and freestanding/hosted runtime? & Compiler command capture, target triplet, parser log.\\
179-guideline mapping & Which source IDs are checked, partially supported, aliased or unsupported? & Versioned check-map with exact semantic scope.\\
System analysis & Are caller/callee, ODR/linkage and whole-program flows understood? & Link-map coverage, unit set and resolution logs.\\
Undecidable rules & Does the engine report possible paths, confidence and assumptions? & Path witness or conservative model narrative.\\
Configuration & Are defines, include paths, templates and code generators reproduced? & compile\_commands.json, macro manifest, preprocessed snapshots.\\
Reliability & Are false-positive and missed-detection rates measured on representative fixtures? & Evaluation corpus and repeatable regression results.\\
Assurance & What level of tool validation is required by project safety/security process? & Qualification argument and issue-tracking process.\\
\bottomrule
\end{longtable}
\normalsize

\section{Analysis scope and coverage}
A Single Translation Unit rule is checkable translation unit by translation unit with a complete and correct engine. A System rule may require call graph, symbol declarations, templates, macro variants and other translation units. Auditors should inspect the linkable program configuration, not only the repository root.

Recommended coverage categories: Verified supported, Partially supported, Manual/Directive, Unsupported, and Not yet evaluated. A checker with a vendor mapping to a rule ID still requires testing with fixture cases and target project code.


\section{Test fixtures for confidence and precision}
\begin{enumerate}
\item Write one minimally violating and one targeted correction specimen for each tool mapping; include source amplification and exception edge cases.
\item For System rules, construct multi-file tests demonstrating that single-TU scans may miss the issue.
\item For undecidable rules, build feasible and infeasible path tests and record uncertainty.
\item Generate false-positive/false-negative examples from accepted production findings and sanitize sensitive code.
\item Version the fixture inputs and analyzer settings alongside tool version and compiler invocation.
\item Treat any change in analyzer behavior as a test failure requiring review, not a reason to disable all relevant diagnostics.
\end{enumerate}

\section{Finding taxonomy}
\small\setlength{\tabcolsep}{3pt}\renewcommand{\arraystretch}{1.15}
\begin{longtable}{p{0.26\linewidth}p{0.62\linewidth}}
\toprule
Disposition & Technical requirement\\ \midrule
\endfirsthead
\toprule
Disposition & Technical requirement\\ \midrule
\endhead
Confirmed violation & Demonstrated violation of the applicable MISRA guideline; enter remediation or controlled deviation path.\\
Compliant by source exception & Quote or cite the specific source-defined exception and why each precondition applies.\\
Analyzer false positive & Analyzer incorrectly reports a rule violation; show a reproducible semantic argument.\\
Undecidable / potential path & Preserve uncertainty and a feasibility assessment; do not mark compliant merely because a test did not trigger it.\\
Out-of-scope code/configuration & Document build and applicability evidence. Do not hide shipping variants.\\
Tool limitation / unsupported & Record residual manual review or a compensating verification method.\\
\bottomrule
\end{longtable}
\normalsize

\section{Automated result normalization}
\begin{lstlisting}[language={},caption={Proposed analysis-record schema (illustrative JSON)}]
{
  "source_edition": "MISRA C++:2023",
  "guideline_id": "Rule 6.8.2",
  "target_config": "embedded-debug-armv7",
  "tool": {"name": "<vendor>", "version": "<version>"},
  "coverage": "verified-supported",
  "analysis_scope": "Single Translation Unit",
  "finding": {"file": "src/example.cpp", "line": 25},
  "triage": "confirmed-violation",
  "review_evidence": "ticket-or-proof-id"
}
\end{lstlisting}

\section{Evidence-based quality indicators}
Track source-rule coverage, complete-build-variant coverage, tool confidence, age of unresolved actual violations, deviation review latency, and recurrence. These are management measures, not alternative MISRA category definitions. Always separately report Mandatory violations regardless of aggregate percent passing.

Recommended pipeline threshold: stop release where Mandatory violations are confirmed or where actual Required violations lack authorized dispositions. A CI gate must account for partial tool coverage; no findings does not prove guideline compliance.


\end{document}
